\SourceLeaf{385}%
[Weiter ergibt sich aus den Gleichungen 16:]
\begin{equation}\tag*{[19]}
\left\{\begin{aligned}
Bc''-Cb''&=\frac{-am'+\alpha m''}{m}\\
Ca''-Ac''&=\frac{-bm'+\beta m''}{m}\\
Ab''-Ba''&=\frac{-cm'+\gamma m''}{m}.
\end{aligned}\right.
\end{equation}
[Durch Vertauschung von $t$ und $u$ erhält man aus 15:
\[
\left\{\begin{aligned}
\alpha''&=\frac{\alpha m''-am'}{m}-A(\alpha A''+\beta B''+\gamma C'')\\
\beta''&=\frac{\beta m''-bm'}{m}-B(\alpha A''+\beta B''+\gamma C'')\\
\gamma''&=\frac{\gamma m''-cm'}{m}-C(\alpha A''+\beta B''+\gamma C'')
\end{aligned}\right.
\]
und leitet hieraus leicht die Relationen ab:]
\begin{equation}\tag*{[20]}
\left\{\begin{aligned}
a'+\alpha''&=-Ammn\\
b'+\beta''&=-Bmmn\\
c'+\gamma''&=-Cmmn.
\end{aligned}\right.
\end{equation}
[Werden nunmehr in den Gleichungen 15 und 16 für $a,b,c$; $\alpha,\beta,\gamma$ die Werthe aus den Gleichungen 12 und 13 eingesetzt, so findet man:]
\begin{equation}\tag*{[21]}
\left\{\begin{aligned}
na'&=\left(\frac{m'}m-AA'\right)(A'+CB''-BC'')-AB'(B'+AC''-CA'')\\
&\quad-AC'(C'+BA''-AB'')-\frac{m''}m(A''+BC'-CB'),\\[4pt]
nb'&=\left(\frac{m'}m-BB'\right)(B'+AC''-CA'')-BC'(C'+BA''-AB'')\\
&\quad-BA'(A'+CB''-BC'')-\frac{m''}m(B''+CA'-AC'),\\[4pt]
nc'&=\left(\frac{m'}m-CC'\right)(C'+BA''-AB'')-CA'(A'+CB''-BC'')\\
&\quad-CB'(B'+AC''-CA'')-\frac{m''}m(C''+AB'-BA')
\end{aligned}\right.
\end{equation}
[und]
\begin{equation}\tag*{[22]}
\left\{\begin{aligned}
na''&=\frac{m'}m(A''+BC'-CB')+\frac{m''}m(A'+CB''-BC'')\\
&\quad-A(A'A''+B'B''+C'C'')\\[4pt]
nb''&=\frac{m'}m(B''+CA'-AC')+\frac{m''}m(B'+AC''-CA'')\\
&\quad-B(A'A''+B'B''+C'C'')\\[4pt]
nc''&=\frac{m'}m(C''+AB'-BA')+\frac{m''}m(C'+BA''-AC'')\\
&\quad-C(A'A''+B'B''+C'C''),
\end{aligned}\right.
\end{equation}
